\section{Standing}
\label{sec:standing}

A credit community must answer one question about each of its members: how much will it
stand behind them. The answer here is a minimum cut --- not a figure the ledger stores,
but one computed from the stakes and the outstanding obligations that it does.

\subsection{Underwriters}

\begin{definition}[Underwriter]
\label{def:underwriter}
An \emph{underwriter} is a member who has declared a \emph{supply}
$\supply(u) \in \mathbb{N}$: the amount of the community's credit they are willing to
stand behind. $\uw$ denotes the set of underwriters. Genesis seeds the first underwriters
directly; thereafter a supply rises only by an endorsed amendment
(\S\ref{sec:seed-amendments}). A member may lower one at any time, never below what they
already stand behind, so restating an unchanged declaration is always admissible.
\end{definition}

Each half of that rule is load-bearing. Seating by ceremony stops a free key declaring a
billion, for a reason no cap drawn from inside the community can supply: a commitment
brought from outside cannot be expressed by any quantity the ledger confers, so the ledger
asks the community instead. The floor at the current declaration keeps a founder ---
whose own capacity is zero until somebody backs them --- from giving up a standing promise
by restating it. An underwriter is not a rank and not a privilege but an accepted
liability: a declaration that if the members this underwriter backs fail, this much of the
loss is theirs (\S\ref{sec:recourse}).

\paragraph{The community is not limited by its founders.} $\uw$ grows as the community
endorses new underwriters, so its total underwriting is not fixed at genesis and does not
die with the people who started it. Founders are simply the first underwriters.

\paragraph{A declaration may not be made against capacity.} The obvious alternative ---
let a member declare up to their own $\Capa$, since a key nobody backs then has nothing
to declare --- reads correctly for one member at a time and issues hollow insurance in
aggregate: capacity is what the community conferred, so such a declaration becomes a
source arc feeding the next member's capacity, which becomes their declaration. Measured,
one founder with a supply of $100$ and twelve joiners wash-backing each other, each
declaring the maximum that rule allows, reach a declared total of $204{,}900$; each then
backs one fresh account, and honest creditors lend those accounts $204{,}800$, every unit
of it carrying the ledger's own \emph{insured} label, with every invariant in
\S\ref{sec:verification} satisfied throughout. That the twelve measured \emph{together}
can owe exactly $100$ is Corollary~\ref{cor:self-uw}, a statement about the twelve; the
credit went to the accounts they back, which are not in that set. So a supply has one
source, and $\sum_{u\in\uw}\supply(u)$ is the community's insured credit rather than an
upper bound on it.

\paragraph{A supply may not be withdrawn out from under live credit.} Any member may
reduce a declaration, but only down to the flow currently committed through them: the
decay floor below, applied to a source arc instead of an edge, because the debt did not
shrink when the underwriter changed their mind. Without it withdrawal breaks
Theorem~\ref{thm:cut} directly --- measured, one of six underwriters leaving a fully drawn
community leaves capacity $12{,}500$ against $15{,}000$ outstanding.

\paragraph{An underwriter's failure is bounded.} Supply bounds not only what an
underwriter may confer on any one member but the total they may confer on everybody, so
the worst an underwriter can do --- by folly, by malice, or by backing a coalition of
identities they control --- is lose their declared supply.

\subsection{The stake graph}

\begin{definition}[Stake]
\label{def:stake}
$\stake(c,d)\in\mathbb{N}$ is the \emph{stake} of creditor $c$ on debtor $d$: directed,
denominated in minor units, updated only by discharge. When $d$ has repaid $R$ of an
obligation of $A$ from $d$ to $c$ --- at every payment, with $R$ the obligation's
cumulative repayment, so $R = A$ when it settles ---
\[
  \stake(c,d) \;\longleftarrow\; \max\!\big(\stake(c,d),\; \min(R,\ \Gamma(c))\big),
\]
where $\Gamma(c)$ is what $c$ may confer: $\supply(c)$ if $c\in\uw$, and $\Capa(c)$
otherwise. Read over the installment rather than over the obligation, a member who honoured
$1{,}000$ in two halves would hold $500$ beside one who paid at once and held $1{,}000$,
for the same evidence; the peak over cumulative repayment equals the rule at settlement
and counts the installments on the way.
\end{definition}

Three properties of Definition~\ref{def:stake} carry the security argument.

\paragraph{It is directed.} That $c$ is willing to carry $d$ says nothing about whether
anyone is willing to carry $c$. Held symmetrically, extending credit would raise the
\emph{lender's} own limit --- a member could grow their capacity by lending, which is
backwards, and which a symmetric graph cannot express its way out of.

\paragraph{It is a peak, not a sum.} Running the same cycle of obligations between one
pair a thousand times raises $\stake$ exactly as high as running it once, so a wash loop
cannot accumulate however patient its operators are.

\paragraph{It is capped by what the creditor may confer.} Nobody can pass on more
standing than they hold. Between two accounts that may confer nothing, every settlement
stakes $\min(A,0)=0$, so no volume of traffic creates an edge at all.

$\Gamma$ is read live at \emph{settlement}, not at acceptance. For a non-underwriter that
is $\Capa(c)$ as Definition~\ref{def:capacity} gives it, net of the credit already
standing on $c$, whose own outstanding insured debt has reserved the very edges the query
runs over --- so a creditor drawn to their own ceiling passes on nothing when a debtor
repays them: what they may confer is what the community currently carries them for, and
they are already using it. The two readings differ only when the creditor's own backing
moved while the obligation was outstanding, and then the present figure is the true one.

\begin{remark}[Why stake and not volume]
\label{rem:not-volume}
Letting $\stake$ record settled volume, the obvious measure of a trading relationship,
is unsound: settlement requires two signatures and no delivery, so volume is free to
fabricate, and two cooperating accounts can sign arbitrarily large settlements having
moved nothing. A stake is not free, because it is capped by a quantity the fabricator
had to earn.
\end{remark}

\begin{remark}[What standing measures]
\label{rem:owed-not-sold}
Because flow runs creditor to debtor, standing is evidence of having \emph{owed and
paid}. Selling earns none of it: measured, $2000$ settled sales to twenty well-backed
members leave the seller's capacity at zero, while one honoured purchase confers standing
at once. Capacity asks whether the community will carry your debt, and the evidence for
that is a debt you carried; producing value is different evidence about a different
question, and the asymmetry is real.

It cannot be relaxed within the model. Crediting the selling side by a fraction $\varphi$
of what the buyer may confer is not a dial: every $\varphi>0$ reaches the same fixed
point, the community ceiling, differing only in rate, because the fraction is applied to
a quantity the fraction itself grows and Definition~\ref{def:stake}'s peak rule makes it a
ratchet. Nor can selling raise some other quantity instead: granting a member whose
capacity is $1000$ the right to confer $10^{6}$ leaves the newcomer they back at $1000$,
and a seller with capacity zero granting $10^{6}$ to each of twenty accounts leaves every
one of them at zero. A flow network holds exactly two things, source arcs and edge
weights; selling can only touch edge weights, and an edge carries only what reaches its
tail. To make selling create anything the seller must become a source, and by
Assumption~\ref{ass:free-signatures} a source is an external commitment. What is left for
a trading history lies outside the flow model: rationing when the ceiling binds, and
evidence for the uninsured tier, where the creditor bears the risk and may weigh whatever
they like.
\end{remark}

\subsection{Capacity}

\begin{definition}[Capacity]
\label{def:capacity}
For a set of accounts $T$, the \emph{gross} capacity $\Capa^{0}(T)$ is the maximum flow
into $T$ over the stake graph, where each underwriter $u \notin T$ supplies its whole
declared supply:
\[
  \Capa^{0}(T) \;=\; \mathrm{maxflow}\Big(\textstyle\sum_{u\in\uw\setminus T}\supply(u)
  \;\longrightarrow\; T\Big).
\]
Let $\res(c,d)$ be flow reserved against outstanding insured credit,
$\free(c,d) = \stake(c,d) - \res(c,d)$, and $\mathrm{com}(u)$ the flow already committed
through underwriter $u$. The \emph{residual} capacity $\Capa(T)$ is the same maximum
taken over what is left:
\[
  \Capa(T) \;=\; \mathrm{maxflow}\Big(\textstyle\sum_{u\in\uw\setminus T}
  \big(\supply(u)-\mathrm{com}(u)\big)
  \;\longrightarrow\; T\Big).
\]
\end{definition}

\begin{remark}[Which one a sentence means]
\label{rem:two-capacities}
The two answer different questions and are not interchangeable. $\Capa(T)$ is
\emph{how much more} $T$ may borrow, so it is what a transaction consults and what a
member is shown, and it reads zero exactly when the ceiling is legitimately full.
$\Capa^{0}(T)$ is \emph{how much the community stands behind} $T$ altogether, so it is
what a bound on outstanding credit must be stated against --- comparing credit already
drawn against $\Capa(T)$ compares a quantity with itself subtracted out. Every claim
below says which it means.
\end{remark}

Both subtractions are one statement made on the two halves of a path, and neither is
optional (Remark~\ref{rem:charge-supply}): $\free$ nets what the stake edges are already
carrying, $\supply-\mathrm{com}$ what the source arc is. A supply arc read at its full
declared value lets every debtor of one underwriter draw the whole supply --- the
difference between a bound and a suggestion.

Underwriters \emph{inside} $T$ supply nothing to it. This clause is what makes a
coalition unable to underwrite itself: appointing your own accomplices as underwriters
adds nothing to what your coalition can borrow, because the measurement excludes them
along with the rest of the coalition. The set form is the security statement and the
singleton $\Capa(x)$ is what a transaction consults; because capacity is a cut, the bound
applies to any coalition as a whole, and splitting across identities gains nothing, since
stakes internal to a set never cross its own boundary.

\subsection{Capacity is generative, not scarce}

Definition~\ref{def:capacity} does not divide a fixed pot, and the difference decides
whether a community can grow. Backing somebody costs the backer nothing: $\Gamma$ caps
what a member \emph{may} confer, it is not a balance they spend, and a member who backs
twenty people has exactly the capacity they started with. Standing also propagates
undiminished --- a member backed at $2500$ may back another at $2500$, and so on outward,
so capacity reaches strangers along chains of ordinary trade. What is bounded is
\emph{simultaneous use}: limits propagate freely, while the credit actually outstanding
at any moment must cross a cut and is therefore bounded by Theorem~\ref{thm:cut}.

A community's working credit is not the sum of its underwriters' supplies. That sum is an
upper bound rather than a measure (Corollary~\ref{cor:uw-bound}): it rises whenever
anybody declares, and a declaration made against standing the community itself conferred
adds nothing the community can carry in aggregate. Working credit rises when somebody
accepts a liability the community was not already carrying --- an underwriter who can pay
from \emph{outside} the set being measured --- and \S\ref{sec:seed-amendments} is the only
door to that.

\subsection{Reservation}

\begin{definition}[Reservation]
\label{def:reservation}
Accepting an insured obligation of $A$ for debtor $d$ computes an augmenting flow of
exactly $A$ into $d$ and adds the flow it used to $\res$. Settlement or discharge
releases it.
\end{definition}

Reservation states that a unit of standing backs one obligation at a time, and it is what
holds the bound across \emph{independent} creditors. Without it each creditor evaluates a
pristine graph and lends against the same edges, and aggregate exposure overshoots the
true cut by the number of creditors --- measured at $200\times$ with two hundred of them.
The overshoot is unbounded, not a constant factor.

\begin{remark}[What reservation must charge]
\label{rem:charge-supply}
Definition~\ref{def:reservation} says \emph{the flow it used}, and the augmenting path
runs from an underwriter's supply arc through stake edges to the debtor. Both are charged,
and an obligation records exactly what it took, so settlement returns exactly that.
Charging only the stake edges leaves the supply arc at its full declared value on the next
query, so two debtors backed by the same underwriter each draw the whole supply and the
overshoot is the number of debtors that underwriter backs --- measured with the supply arc
uncharged, one supply of $2500$ carries $250{,}000$ across a hundred of them. Returning
less than was taken is the same error reversed: a release proportional across the arcs
incident to the debtor strands the upstream half of every multi-hop path, so settling an
obligation destroys capacity nothing was behind. Exactness also disposes of the question
of release \emph{order}, since there is nothing left to choose.
\end{remark}

\begin{remark}[Why this needs a total order]
\label{rem:order}
The \emph{total} drawn against a region is order-invariant, but the \emph{allocation} is
not: capacity is first-come-first-served, so which creditor lends depends on the
sequence in which acceptances apply. A conserved budget over a shared graph requires a
serialization point --- the argument that forces a total order for double-spend
prevention, arriving by a different route.
\end{remark}

\subsection{Decay}

Stakes decay per epoch by the governed ratio $\kConstStakeDecayNum/\kConstDecayDen$, so
standing reflects present backing rather than history. An edge \textbf{never decays
below its live reservation}: the debt did not shrink because the evidence aged, so the
collateral behind it may not either. Free capacity absorbs decay first and is floored at
zero, and an edge wholly committed to an outstanding obligation does not decay at all
until that obligation clears.

In the unit a member plans in: at the genesis value an unrenewed edge halves in about
thirty epochs, keeps $12\%$ after a quarter, $1.5\%$ after half a year and $0.02\%$ after
a year; integer decay drops at least one minor unit an epoch below $43$, so an edge of
$500$ is gone in $346$. The floor protects only what is still owed, so decay is paid by
the member who has repaid and not by the member who is drawn. A community trading on a
seasonal rhythm --- a harvest a year --- is erased between seasons at this ratio, and
sizes the governed constant before it sizes anything else (\S\ref{sec:seedsizing}): at the
top of its range a year keeps $69\%$.

\begin{remark}[A community that stops trading loses its credit]
\label{rem:winddown}
A community whose members cease to renew their backing watches its capacity fall to zero
--- at the genesis ratio, to two ten-thousandths in a year. A community that has stopped
trading has no credit, and a system reporting otherwise would be reporting a relationship
nobody currently has. What must not happen is that the \emph{ability} to renew dies out,
which is why $\uw$ is open to any member rather than fixed at genesis.
\end{remark}

\subsection{Insured and uninsured credit}

Capacity bounds what the community \emph{underwrites}, not what a member may choose to
risk. Within the debtor's capacity an obligation is \textbf{insured}: it reserves flow,
and the recourse machinery of \S\ref{sec:recourse} stands behind it. Beyond it a creditor
may still lend, and the obligation is \textbf{uninsured}: it reserves nothing, triggers no
community recourse on default, and the creditor bears it alone.

This dissolves the bootstrap. Were capacity a permission, a community whose accounts all
begin at zero would deadlock, since the first trade could never happen; instead first
trades are uninsured, they settle, they create stakes, and capacity is their residue.
Uninsured credit still stakes under Definition~\ref{def:stake}, so an uninsured loan from
an account that may confer nothing confers nothing.

Capacity then grows \emph{insured across creditors and uninsured within one}. The cut sums
every honoured relationship, and a loan from any creditor is insured against the whole of
it, so each settlement that writes a new creditor's edge raises what the next loan may be
insured for. Measured: one uninsured loan of $50$, then eight loans inside capacity from
three creditors in rotation, take a member through $100$, $200$, $350$, $650$, $1{,}200$,
$2{,}000$ and $2{,}650$ to $3{,}000$ --- the seed's whole reach through those three ---
with $7{,}200$ lent insured on $50$ lent uninsured. Within a single relationship the stake
is a peak, so a limit at one creditor rises only by a loan above it, and reservation being
whole, that loan is uninsured in full: the marginal risk sits with the party holding the
information about the debtor, and the community insures what has been demonstrated.

\begin{remark}[The friction this leaves]
\label{rem:friction}
An unconditional starter allowance would make onboarding frictionless, and is exactly
the hole that makes identities worth minting: $N$ free keys would carry $N$ allowances.
Refusing it moves a real cost onto the newcomer --- somebody must take a first, uninsured
risk on them, and no mechanism here supplies that willingness. That is how a trading
relationship has always begun, a shopkeeper letting a new neighbour run a tab, and this
design's part is to make the risk small, bounded and legible: the amount is the
creditor's own choice; the exposure is exactly that amount and never more
(Cor.~\ref{cor:linear}); the client states at the point of decision that the obligation
is uninsured (\S\ref{sec:medium}); and the bond that same counterparty pays is what seats
the account (\S\ref{sec:implementation}). The risk is what writes the stake, and a
mechanism that removed it would remove what capacity is made of.
\end{remark}

\subsection{What an account is}

An account is a set of keys and the bookkeeping that hangs off them. It carries no tier,
no attestation, no sponsor and no score, and in particular no \emph{admission} status:
there is no rung to be on and none to be promoted from. The status it does carry is the
governance sanction and the member's own departure (\S\ref{sec:model}), and neither is
consulted by Definition~\ref{def:capacity} --- the graph answers the same for a suspended
account as for an active one, and what suspension revokes is origination, not standing.

It is seated by its first \emph{bonded trade} --- an acceptance or a sale naming its key
as a party, signed by that key --- and requires no approval, because creation grants
nothing: an account with no incident stakes has capacity zero by
Definition~\ref{def:capacity} and confers nothing by Definition~\ref{def:stake}. Seating
by trade rather than by a creation transition is a cost argument, not a permission one
(\S\ref{sec:implementation}): a row that only ever appears beside a write somebody with
standing paid for is bounded by that standing, while an unbillable creation is bounded by
nothing and cannot be given a quota that is not also a censorship lever. The counterparty
who carries that first bond is the one already carrying the first uninsured risk
(\S\ref{sec:medium}); nobody else is asked, and nobody may refuse on their behalf. The
row is written \emph{last}, after every check that could refuse the trade, so a refused
transaction leaves nothing behind.

The row is a stock, and its price is one: the seat holds a bond unit of the counterparty's
reach on a layer credit never reads, for as long as the row exists
(Property~\ref{prop:seat}). A bond is reserved and returned, so it prices how fast rows
appear and not how many; no transition returns a seat, because no transition retires a
row. What does retire one is the epoch sweep, a year after the seating, when the row holds
nothing, owes nothing and is named by nothing --- no stake in either direction, no
obligation in any status, no bond, default, supply, vote, pending rotation, panel,
proposal or guardian roll, and no live row it sponsored --- and then the seat and the key
come back. Each of those is something only the member, a counterparty or a ceremony can
put there, so the rule is not a lever anybody can pull against another, and a row that is
empty is not a stock. The last condition is the sweep's own: a seat names its sponsor for
the life of the row it bought, so a sponsor whose own trade has gone quiet stays until the
rows it seated are gone, and goes at the boundary after the last of them.

What an account accumulates is not a property of itself but of the community around it.
Standing is held in the stakes others place on it; no account owns its own capacity, and
none can raise it alone.
